\section{Bring-Up Problems and Changes}
\label{sec:res-failures}

Bring-up required the circuit, sensing, and firmware to be corrected together.
Blocking ADC reads delayed feedback, the initial PA2 interpretation used the
wrong current for output power, and unloaded operation needed a tighter
energy limit. Those problems shaped the final control structure.

\begin{table}[htbp]
\centering
\caption{Bring-up problems, corrections, and design lessons.}
\label{tab:failure-log}
\small
\begin{tabularx}{\textwidth}{@{}L{0.23\textwidth}L{0.36\textwidth}X@{}}
\toprule
Issue & Corrective action & Resulting lesson\\
\midrule
Blocking feedback acquisition & Adopted circular ADC DMA & Acquisition timing must remain compatible with the live control loop\\
PA2 interpreted as output current & Removed PA2-based output-power control; retained PA1 voltage regulation & Match feedback to the physical quantity used by the controller\\
Irregular no-load pulse structure & Revised minimum-cycle handling and no-load duty logic & Schedule complete carrier cycles and limit residual energy explicitly\\
Energy after load removal & Added no-load guarding, voltage windows, and sag recovery paths & Loaded and unloaded operation require distinct control behavior\\
Ambiguous output detail fields & Kept DC, AC detail, and acquisition statistics separate & A range of average readings is not instantaneous output ripple\\
Disagreeing power instruments & Kept the original readings and calculated initial efficiency estimates & State the measurement boundary and synchronization before interpreting losses\\
High-load temperature rise & Recorded camera maxima and the stated test conditions & Follow component temperatures until they settle at the chosen load\\
\bottomrule
\end{tabularx}
\end{table}
\FloatBarrier

PA1 voltage regulation replaced the incorrect power feedback, while PA2
was retained for load classification. Complete-cycle scheduling addressed
irregular pulses, and no-load guarding limited charging after load removal.
The highest calculated V1.0 output was \SI{79.89}{\watt}; further loading toward
the \SI{100}{\watt} objective requires temperature histories alongside the
electrical readings.

The remaining tests must pair input and output power, capture ripple with
known coupling and bandwidth, follow identified temperatures to stabilization,
and record controller state during load and protection events.
